\documentclass[11pt]{article}
\usepackage[margin=1in]{geometry}
\usepackage{amsmath,amssymb}
\usepackage{enumitem}
\title{COMP3242 Week 2 Practice Exam (Questions)}
\date{}
\begin{document}
\maketitle

\section*{Instructions}
Answer all questions. Show derivation steps where needed.

\section*{Q1: Logistic Regression (Calculation)}
For $x=[1,2]^T$, $a=[0.5,-1]^T$, and $b=0.2$:
\begin{enumerate}[label=(\alph*)]
  \item Compute $z=a^Tx+b$.
  \item Compute $\sigma(z)=\frac{1}{1+e^{-z}}$.
  \item If the true label is $y=1$, compute binary cross-entropy loss 
  $L=-\log(\sigma(z))$.
\end{enumerate}

\section*{Q2: XOR and Expressivity (Conceptual)}
Explain why a single linear classifier cannot solve XOR. Then explain how a two-layer perceptron can solve it at a high level.

\section*{Q3: MLP Parameter Count (Calculation)}
A two-layer MLP has input dimension $n=5$, hidden dimension $m=8$, and output dimension $k=3$.
The model is
\[
z=\phi(Ax+b),\quad y=Cz+d,
\]
with $A\in\mathbb{R}^{m\times n}$, $b\in\mathbb{R}^{m}$, $C\in\mathbb{R}^{k\times m}$, $d\in\mathbb{R}^{k}$.
Compute the total number of trainable parameters.

\section*{Q4: Softmax vs Sigmoid (Concept + Application)}
State one appropriate use case for each:
\begin{enumerate}[label=(\alph*)]
  \item Softmax + CrossEntropyLoss
  \item Sigmoid + BinaryCrossEntropy (or BCEWithLogitsLoss)
\end{enumerate}
Also explain in one sentence why they are different.

\section*{Q5: Backprop Shape Tracing (Code/Debug)}
In PyTorch, suppose:
\begin{verbatim}
x:      (64, 20)
A:      (50, 20)
b:      (50,)
C:      (10, 50)
d:      (10,)
\end{verbatim}
For forward pass:
\begin{verbatim}
z = relu(x @ A.T + b)
logits = z @ C.T + d
\end{verbatim}
Give shapes of `z` and `logits`. Then state the expected label shape for `CrossEntropyLoss`.

\end{document}
